\subsection{The Hausdorff completion of a semi-local ring}

PROPOSITION 17. Let A be a commutative ring and \((\mathfrak{m}_{\lambda})_{\lambda \in L}\) a family of ideals of A, distinct from A, such that \(\mathfrak{m}_{\lambda}\) and \(\mathfrak{m}_{\mu}\) are relatively prime for \(\lambda \neq \mu\) . For every family \(s = (s(\lambda))_{\lambda \in L}\) of integers \(\geqslant 0\) , of finite support, set \(\mathfrak{a}_{s} = \bigcap_{\lambda \in L} \mathfrak{m}_{\lambda}^{s(\lambda)}\) (equal to the product of the \(\mathfrak{m}_{\lambda}^{s(\lambda)}\) for the \(\lambda\) such that \(s(\lambda) \neq 0\) ; cf.~Chapter II, § 1, no. 2, Propositions 3 and 5); the \(\mathfrak{a}_{s}\) form a fundamental system of neighbourhoods of 0 with respect to a topology \(\mathcal{T}\) compatible with the ring structure on A; let \(\hat{\mathbf{A}}\) be the Hausdorff completion of A with respect to this topology. On the other hand, for all \(\lambda \in L\) , let \(\mathbf{A}_{\lambda}\) be the ring A with the \(\mathfrak{m}_{\lambda}\) -adic topology and let \(\hat{\mathbf{A}}_{\lambda}\) be its Hausdorff completion. If \(u: \mathbf{A} \to \prod_{\lambda \in L} \mathbf{A}_{\lambda}\) denotes the diagonal homomorphism, \(u\) is continuous and the corresponding homomorphism \(\hat{u}\) :

\[
\hat {\mathbf {A}} \rightarrow \left(\prod_ {\lambda \in \mathbf {L}} \mathbf {A} _ {\lambda}\right) ^ {\wedge} = \prod_ {\lambda \in \mathbf {L}} \hat {\mathbf {A}} _ {\lambda}
\]

(General Topology, Chapter III, § 6, no. 5 and Chapter II, § 3, no. 9, Corollary 2 to Proposition 18) is a topological ring isomorphism.

The first assertion follows from General Topology, Chapter III, § 6, no. 3, Example 3. Let us set \(\mathbf{B} = \prod_{\lambda \in \mathbf{L}} \mathbf{A}_{\lambda}\) ; as the topology on \(\mathbf{A}\) is finer than each of the \(m_{\lambda}\) -adic topologies, the mappings \(\mathrm{pr}_{\lambda} \circ u\) are continuous and hence \(u\) is continuous. Also, \(u(a_s)\) is the intersection of the diagonal \(\Delta\) of \(\mathbf{B}\) and the open set \(\bigcap_{\lambda \in \mathbf{L}} \mathrm{pr}_{\lambda}^{-1}(m_{\lambda}^{s(\lambda)})\) of \(\mathbf{B}\) ; it follows that \(u\) is a strict morphism from the additive group \(\mathbf{A}\) to \(\mathbf{B}\) with image \(\Delta\) . Now \(\Delta\) is dense in \(\mathbf{B}\) . For let \(b = (a_{\lambda})_{\lambda \in \mathbf{L}}\) be an element of \(\mathbf{B}\) ; every neighbourhood of \(b\) in \(\mathbf{B}\) contains a set of the form \(b + \mathbf{V}\) , where \(\mathbf{V} = \bigcap_{\lambda \in \mathbf{L}} \mathrm{pr}_{\lambda}^{-1}(m_{\lambda}^{s(\lambda)})\) for a family \(s = (s(\lambda))_{\lambda \in \mathbf{L}}\) with finite support of integers \(\geqslant 0\) . As the \(m_{\lambda}^{s(\lambda)}\) are relatively prime in pairs (Chapter II, § 1, no. 2, Proposition 3), there exists \(x \in \mathbf{A}\) such that \(x \equiv a_{\lambda}\) (mod. \(m_{\lambda}^{s(\lambda)}\) ) for all \(\lambda\) (loc. cit., Proposition 5) and hence \((b + \mathbf{V}) \cap \Delta \neq \varnothing\) . The Hausdorff completion of the group \(\mathbf{B}/\Delta\) is then \(\{0\}\) ; applying Lemma 2 of no. 12 to the exact sequences \(0 \to \mathbf{A} \xrightarrow{u} \mathbf{B}, \mathbf{A} \xrightarrow{u} \mathbf{B} \to \mathbf{B}/\Delta\) , we see that \(\hat{u}\) is an isomorphism of \(\hat{\mathbf{A}}\) onto \(\hat{\mathbf{B}}\) .

COROLLARY. Let A be a principal ideal domain and P a representative system of extremal elements of A (Algebra, Chapter VII, § 1, no. 3). The topology on A with respect to which the ideals \(\neq0\) of A form a fundamental system of neighbourhoods of 0, which is compatible with the ring structure on A, is Hausdorff and the completion of A with this topology is canonically isomorphic to the product of the completions of A with respect to the \(\pi\) -adic topologies, where \(\pi\) runs through P.

The principal ideals \((\pi)\) where \(\pi \in P\) are maximal and distinct and hence relatively prime, we have already seen (no. 12, Example 3) that the \(\pi\) -adic topologies are Hausdorff and hence so is the topology defined in the statement of Proposition 17, which is finer than each of the \(\pi\) -adic topologies.

If the Corollary to Proposition 17 is applied when A = Z, we denote by \(\hat{Z}\) the completion of Z with respect to the topology for which all the ideals \(\neq0\) of Z form a fundamental system of neighbourhoods of 0, the ring isomorphic to the product \(\prod_{p\in P}Z_{p}\) of the rings of p-adic integers (P being the set of prime numbers).

\textbf{Remarks}\par

\begin{enumerate}
\def\labelenumi{(\arabic{enumi})}
\item
  Clearly, under the conditions of Proposition 17, the topology \(\mathcal{T}\) is the least upper bound of the \(m_{\lambda}\) -adic topologies on \(\mathbf{A}\) .
\item
  Every closed ideal \(\mathfrak{a}\) of \(\prod_{\lambda \in L} \hat{\mathbf{A}}_{\lambda}\) is identical with the product of its projections \(\mathfrak{a}_{\lambda} = \operatorname{pr}_{\lambda}(\mathfrak{a})\) , which are closed ideals in the \(\hat{\mathbf{A}}_{\lambda}\) ; for \(\hat{\mathbf{A}}_{\lambda}\) is canonically identified with a closed ideal \(\mathbf{A}_{\lambda}'\) of \(\prod_{\lambda} \hat{\mathbf{A}}_{\lambda}\) and \(\mathfrak{a}_{\lambda}\) with \(\mathfrak{a} \cap \mathbf{A}_{\lambda}'\) (Algebra, Chapter I, § 8, no. 10, Proposition 6), the sum of the \(\mathfrak{a}_{\lambda}\) is dense in the product \(\prod_{\lambda} \mathfrak{a}_{\lambda}\) (General Topology, Chapter III, § 2, no. 9, Proposition 25) and the latter is closed in \(\prod_{\lambda} \hat{\mathbf{A}}_{\lambda}\) , whence our assertion.
\end{enumerate}

PROPOSITION 18. Let A be a commutative ring, \((\mathfrak{m}_i)_{1 \leqslant i \leqslant q}\) a finite family of distinct maximal ideals of A, r the product ideal \(\mathfrak{m}_1\mathfrak{m}_2 \ldots \mathfrak{m}_q = \mathfrak{m}_1 \cap \mathfrak{m}_2 \cap \cdots \cap \mathfrak{m}_q\) and S the multiplicative subset \(\bigcap_{i=1}^{q} (\mathbf{A} - \mathfrak{m}_i)\) . Let A be given the r-adic topology, the ring \(\mathbf{B} = \mathbf{S}^{-1}\mathbf{A}\) the rB-adic topology and each of the local rings \(\mathbf{A}_{\mathfrak{m}_i}\) the \((\mathfrak{m}_i\mathbf{A}_{\mathfrak{m}_i})\) -adic topology. Let \(u: \mathbf{A} \to \mathbf{B}\) , \(v_i: \mathbf{B} \to \mathbf{A}_{\mathfrak{m}_i}\) be the canonical homomorphisms (Chapter II, § 2, no.1, Corollary 2 to Proposition 2) and \(v\) the homomorphism \((v_i): \mathbf{B} \to \prod_{i=1}^{q} \mathbf{A}_{\mathfrak{m}_i}\) . The homomorphisms \(u\) and \(v\) are continuous and the corresponding homomorphisms \(\hat{u}: \hat{\mathbf{A}} \to \hat{\mathbf{B}}\) and \(\hat{v}: \hat{\mathbf{B}} \to \prod_{i=1}^{q} (\mathbf{A}_{\mathfrak{m}_i})^{\wedge}\) are topological ring isomorphisms.

\(m_{i} \cap S = \varnothing\) for \(1 \leqslant i \leqslant q\) , hence the ideal \(m_{i}' = m_{i}B\) of B is maximal (Chapter II, §2, no. 5, Proposition 11) and

\[
\mathfrak {r B} = \mathfrak {m} _ {1} ^ {\prime} \cap \mathfrak {m} _ {2} ^ {\prime} \cap \dots \cap \mathfrak {m} _ {q} ^ {\prime}
\]

(Chapter II, §2, no. 4); finally, \(\mathbf{B}_{\mathfrak{m}_i'} = \mathbf{A}_{\mathfrak{m}_i}\) up to a canonical isomorphism (Chapter II, §2, no. 5, Proposition 11). As \(\bar{u}^{-1}(\mathfrak{rB}) = \mathfrak{r}\) and \(\bar{v}_i^{-1}(m_i\mathbf{A}_{\mathfrak{m}_i}) \supset \mathfrak{rB}\) , \(u\) and \(v\) are continuous. Then it is sufficient to prove that,

\[
w = v \circ u: \mathrm{A} \rightarrow \prod_ {i = 1} ^ {q} \mathrm{A} _ {\mathfrak {m} _ {i}},
\]

\(\hat{w}\) is an isomorphism of \(\hat{\mathbf{A}}\) onto \(\prod_{i=1}^{q} \hat{\mathbf{A}}_{m_i}\) , for this result applied to B and the \(m_i'\) will show that \(\hat{v}\) is an isomorphism and therefore also \(\hat{u}\) . Note that every product of powers of the \(m_i\) contains a power of r and hence the r-adic topology is the least upper bound of the \(m_i\) -adic topologies; moreover, if \(A_i\) denotes the

\[
\mathfrak {m} _ {i} \text {-adic}
\]

\[
\phi : \mathrm{A} \rightarrow \prod_ {i = 1} ^ {q} \mathrm{A} _ {i}
\]

\(\hat{\phi}:\hat{\mathbf{A}}\to \prod_{i = 1}^{q}\hat{\mathbf{A}}_i\) is an isomorphism (Proposition 17). Then it all amounts to proving that, if \(u_{i}:\mathbf{A}_{i}\rightarrow \mathbf{A}_{m_{i}}\) is the canonical mapping, \(\hat{u}_i\colon \hat{\mathbf{A}}_i\to \hat{\mathbf{A}}_{m_i}\) is an isomorphism. Now, for all \(n\) , the mapping

\[
u _ {i, n} \colon \mathrm{A} / \mathfrak {m} _ {i} ^ {n} \rightarrow \mathrm{A} _ {\mathfrak {m} _ {i}} / \mathfrak {m} _ {i} ^ {n} \mathrm{A} _ {\mathfrak {m} _ {i}}
\]

derived from \(u_{i}\) by taking quotients is an isomorphism (Chapter II, § 3, no. 3, Proposition 9); our assertion follows from the fact that \(\hat{\mathbf{A}}_i\) (resp. \(\hat{\mathbf{A}}_{\mathfrak{m}_i}\) ) is the inverse limit of the discrete rings \(\mathbf{A} / \mathfrak{m}_i^n\) (resp. \(\mathbf{A}_{\mathfrak{m}_i} / \mathfrak{m}_i^n \mathbf{A}_{\mathfrak{m}_i}\) ) (cf.~no. 6).

Remark (3). We see that an integral domain A can be such that its Hausdorff completion \(\hat{A}\) admits non-zero divisors of zero.

PROPOSITION 19. Let A be a commutative ring and m a maximal ideal of A. The Hausdorff completion \(\hat{\mathbf{A}}\) of A with respect to the m-adic topology is a local ring whose maximal ideal is \(\hat{\mathfrak{m}}\) .

If \(\mathfrak{a} = \bigcap_{k\geqslant 1}\mathfrak{m}^k\) , \(\hat{\mathbf{A}}\) is the completion of the Hausdorff ring \(\mathbf{A} / \mathfrak{a}\) associated with \(\mathbf{A}\) and, as \(\mathfrak{m} / \mathfrak{a}\) is maximal in \(\mathbf{A} / \mathfrak{a}\) , we may assume that \(\mathbf{A}\) is Hausdorff with respect to the \(\mathfrak{m}\) -adic topology. As \(\mathbf{A} / \mathfrak{m}\) and \(\hat{\mathbf{A}} /\hat{\mathfrak{m}}\) are isomorphic rings (no. 12, formula (21)), \(\hat{\mathfrak{m}}\) is maximal in \(\hat{\mathbf{A}}\) . As the topology on \(\hat{\mathbf{A}}\) is defined by the filtration \((\mathfrak{m}^n)^{\wedge}\) (no. 12), the proposition will be a consequence of the following lemma:

LEMMA 3. Let A be a complete Hausdorff topological ring, in which there exists a fundamental system \(\mathfrak{S}\) of neighbourhoods of 0 consisting of additive subgroups of A.

\begin{enumerate}
\def\labelenumi{(\roman{enumi})}
\item
  For all \(x \in \mathbf{A}\) such that \(\lim_{n \to \infty} x^n = 0\) , \(1 - x\) is invertible in \(\mathbf{A}\) and its inverse is equal to \(\sum_{n=0}^{\infty} x^n\) .
\item
  Let \(\mathfrak{a}\) be a two-sided ideal of A such that \(\lim x^n = 0\) for all \(x \in \mathfrak{a}\) . For an element y of A to be invertible, it is necessary and sufficient that its class mod. a be invertible in A/a; in particular a is contained in the Jacobson radical of A.
\end{enumerate}

(i) As

\[
(1 - x) (1 + x + \dots + x ^ {n}) = (1 + x + \dots + x ^ {n}) (1 - x) = 1 - x ^ {n + 1},
\]

it all amounts to proving that the series with general term \(x^{n}\) is convergent in A; now, by hypothesis, for every neighbourhood \(V \in \mathfrak{S}\) of 0 in A, there exists an integer p \textgreater{} 0 such that \(x^{n} \in V\) for all \(n \geqslant p\) . We conclude that

\[
x ^ {p} + x ^ {p + 1} + \dots + x ^ {q} \in \mathbf {V}
\]

for all \(q \geqslant p\) and our assertion then follows from Cauchy's criterion (General Topology, Chapter III, § 5, no. 2, Theorem 1).

\begin{enumerate}
\def\labelenumi{(\roman{enumi})}
\setcounter{enumi}{1}
\tightlist
\item
  Suppose that there exists \(y' \in A\) such that \(yy' \equiv 1 \pmod{a}\) and \(y'y \equiv 1 \pmod{a}\) . The hypothesis on \(a\) implies, by (i), that \(yy'\) and \(y'y\) are invertible in \(A\) and hence \(y\) is invertible in \(A\) . In particular, every \(x \in a\) is such that \(1 - x\) is invertible in \(A\) and, as \(a\) is a two-sided ideal of \(A\) , it is contained in the Jacobson radical of \(A\) (Algebra, Chapter VIII, § 6, no. 3, Theorem 1).
\end{enumerate}

Having established this lemma, it is sufficient to apply it to the topological ring \(\hat{\mathbf{A}}\) and the ideal \(\hat{\mathfrak{m}}\) , as, for all \(x \in \hat{\mathfrak{m}}\) , \(x^n \in (\hat{\mathfrak{m}})^n \subset (\mathfrak{m}^n)^\wedge\) and the sequence \((x^n)\) therefore tends to 0.

If we take A = Z, every maximal ideal of Z is of the form pZ where p is prime. The ring of p-adic numbers \(Z_{p}\) is then a local ring of which \(pZ_{p}\) is the maximal ideal (Corollary 2 to Proposition 16) and whose residue field is isomorphic to \(Z/pZ = F_{p}\) , and \(\mathbf{Z}_{(p)}\) with the \(pZ_{(p)}\) -adic topology is identified with a topological subring of \(Z_{p}\) containing Z.

COROLLARY. Let A be a semi-local ring (Chapter II, § 3, no. 5), \(\mathfrak{m}_i\) its distinct maximal ideals ( \(1 \leqslant i \leqslant q\) ) and

\[
\mathfrak {r} = \mathfrak {m} _ {1} \cap \mathfrak {m} _ {2} \cap \dots \cap \mathfrak {m} _ {q}
\]

its Jacobson radical. The Hausdorff completion \(\hat{\mathbf{A}}\) of \(\mathbf{A}\) with respect to the \(\mathfrak{r}\) -adic topology is a semi-local ring, canonically isomorphic to the product \(\prod_{i=1}^{q} \hat{\mathbf{A}}_{\mathfrak{m}_i}\) , where \(\hat{\mathbf{A}}_{\mathfrak{m}_i}\) is the Hausdorff completion ring of the local ring \(\mathbf{A}_{\mathfrak{m}_i}\) with respect to the \((\mathfrak{m}_i\mathbf{A}_{\mathfrak{m}_i})\) -adic topology.

